% GENERATED FILE. Edit canonical lessons, then run npm run generate:books.
% canonical-source-hash: fnv1a64:47ed0189f6431c67
% canonical-lessons: TA-C317-mutanmuraiyaka, TA-C317-tirumpattirumpa, TA-C317-totarcciyaka, TA-C317-utanukkutan, TA-C317-atuttatuttu, TA-R317-review-310-317

\chapter{Again and Again}
\label{ch:ta-again-and-again}
\hlchaptermodality{317}

\begin{chapteropening}
\textbf{By the end of this chapter:} \emph{I can say how often and how soon things happened.}

Complete the last lesson of chapter 317: I can say how often and how soon things happened.
\end{chapteropening}

\section[\texorpdfstring{\ta{முதன்முறையாக} (mutaṉmuṟaiyāka)}{முதன்முறையாக (mutaṉmuṟaiyāka)}]{\ta{முதன்முறையாக} (mutaṉmuṟaiyāka) — for the first time}

\label{lesson:TA-C317-mutanmuraiyaka}



\begin{quote}
Before the new one: say the Tamil for after that, then the Tamil for at the same time.
\end{quote}

\subsection*{You'll want to know: \ta{முதன்முறையாக}}
\textbf{\ta{முதன்முறையாக}} — \emph{mutaṉmuṟaiyāka} — \textquotedblleft{}for the first time\textquotedblright{}.

\begin{quote}
\ta{நான்} \ta{முதன்முறையாக} \ta{விமானத்தில்} \ta{பயணம்} \ta{செய்தேன்}.
\end{quote}

\begin{quote}
\emph{nāṉ mutaṉmuṟaiyāka vimāṉattil payaṇam ceytēṉ.}
\end{quote}

\begin{quote}
I travelled by plane for the first time.
\end{quote}

\ta{முதல்} + \ta{முறை}; \ta{ல்} turns to \ta{ன்} before \ta{ம}.

\begin{cousinweb}
\textbf{\ta{முதன்முறையாக}} (\emph{mutaṉmuṟaiyāka}) is built from native \ta{முதல்}, \textquotedblleft{}first\textquotedblright{} (DEDR 4950), and \ta{முறை}, \textquotedblleft{}order, turn, time\textquotedblright{} (DEDR 5015), with the adverb ending -\ta{ஆக}.
\end{cousinweb}

\begin{grammarlens}[title={What you've built}]
One more word for telling what happened.
\end{grammarlens}

\subsection*{Your turn}
\begin{itemize}
\raggedright
  \item \emph{Say it:} \emph{mutaṉmuṟaiyāka}
  \item \emph{Say it:} \emph{mutaṉmuṟaiyāka}, once more
  \item \emph{Say it:} \emph{nāṉ mutaṉmuṟaiyāka vimāṉattil payaṇam ceytēṉ.} — I travelled by plane for the first time.
  \item \emph{Recall:} say the Tamil for after that, then the Tamil for at the same time, then say \emph{mutaṉmuṟaiyāka} again
\end{itemize}

\begin{tcolorbox}[breakable,colback=teal!4,colframe=teal!35!black,title={Before you move on}]
What does \ta{முதன்முறையாக} mean? (\textquotedblleft{}For the first time\textquotedblright{}.) Say it once more.
\end{tcolorbox}
\section[\texorpdfstring{\ta{திரும்பத்திரும்ப} (tirumpattirumpa)}{திரும்பத்திரும்ப (tirumpattirumpa)}]{\ta{திரும்பத்திரும்ப} (tirumpattirumpa) — again and again}

\label{lesson:TA-C317-tirumpattirumpa}



\begin{quote}
Before the new one: say the Tamil for at the same time, then the Tamil for for the first time.
\end{quote}

\subsection*{You'll want to know: \ta{திரும்பத்திரும்ப}}
\textbf{\ta{திரும்பத்திரும்ப}} — \emph{tirumpattirumpa} — \textquotedblleft{}again and again\textquotedblright{}.

\begin{quote}
\ta{அவன்} \ta{அதே} \ta{கேள்வியைத்} \ta{திரும்பத்திரும்பக்} \ta{கேட்டான்}.
\end{quote}

\begin{quote}
\emph{avaṉ atē kēḷviyait tirumpattirumpak kēṭṭāṉ.}
\end{quote}

\begin{quote}
He asked the same question again and again.
\end{quote}

The infinitive of \ta{திரும்பு}, doubled.

\begin{cousinweb}
\textbf{\ta{திரும்பத்திரும்ப}} (\emph{tirumpattirumpa}) doubles the infinitive of native \ta{திரும்பு}, \textquotedblleft{}to turn, turn back\textquotedblright{} (DEDR 3246).
\end{cousinweb}

\begin{grammarlens}[title={What you've built}]
One more word for telling what happened.
\end{grammarlens}

\subsection*{Your turn}
\begin{itemize}
\raggedright
  \item \emph{Say it:} \emph{tirumpattirumpa}
  \item \emph{Say it:} \emph{tirumpattirumpa}, once more
  \item \emph{Say it:} \emph{avaṉ atē kēḷviyait tirumpattirumpak kēṭṭāṉ.} — He asked the same question again and again.
  \item \emph{Recall:} say the Tamil for at the same time, then the Tamil for for the first time, then say \emph{tirumpattirumpa} again
\end{itemize}

\begin{tcolorbox}[breakable,colback=teal!4,colframe=teal!35!black,title={Before you move on}]
What does \ta{திரும்பத்திரும்ப} mean? (\textquotedblleft{}Again and again\textquotedblright{}.) Say it once more.
\end{tcolorbox}
\section[\texorpdfstring{\ta{தொடர்ச்சியாக} (toṭarcciyāka)}{தொடர்ச்சியாக (toṭarcciyāka)}]{\ta{தொடர்ச்சியாக} (toṭarcciyāka) — continuously}

\label{lesson:TA-C317-totarcciyaka}



\begin{quote}
Before the new one: say the Tamil for for the first time, then the Tamil for again and again.
\end{quote}

\subsection*{You'll want to know: \ta{தொடர்ச்சியாக}}
\textbf{\ta{தொடர்ச்சியாக}} — \emph{toṭarcciyāka} — \textquotedblleft{}continuously\textquotedblright{}.

\begin{quote}
\ta{மூன்று} \ta{நாள்} \ta{தொடர்ச்சியாக} \ta{மழை} \ta{பெய்தது}.
\end{quote}

\begin{quote}
\emph{mūṉṟu nāḷ toṭarcciyāka maḻai peytatu.}
\end{quote}

\begin{quote}
It rained for three days without a break.
\end{quote}

From \ta{தொடர்}, to continue.

\begin{cousinweb}
\textbf{\ta{தொடர்ச்சியாக}} (\emph{toṭarcciyāka}) is built on \ta{தொடர்ச்சி}, \textquotedblleft{}continuation\textquotedblright{}, from native \ta{தொடர்}, \textquotedblleft{}to follow uninterruptedly\textquotedblright{} (DEDR 3480), with -\ta{ஆக}.
\end{cousinweb}

\begin{grammarlens}[title={What you've built}]
One more word for telling what happened.
\end{grammarlens}

\subsection*{Your turn}
\begin{itemize}
\raggedright
  \item \emph{Say it:} \emph{toṭarcciyāka}
  \item \emph{Say it:} \emph{toṭarcciyāka}, once more
  \item \emph{Say it:} \emph{mūṉṟu nāḷ toṭarcciyāka maḻai peytatu.} — It rained for three days without a break.
  \item \emph{Recall:} say the Tamil for for the first time, then the Tamil for again and again, then say \emph{toṭarcciyāka} again
\end{itemize}

\begin{tcolorbox}[breakable,colback=teal!4,colframe=teal!35!black,title={Before you move on}]
What does \ta{தொடர்ச்சியாக} mean? (\textquotedblleft{}Continuously\textquotedblright{}.) Say it once more.
\end{tcolorbox}
\section[\texorpdfstring{\ta{உடனுக்குடன்} (uṭaṉukkuṭaṉ)}{உடனுக்குடன் (uṭaṉukkuṭaṉ)}]{\ta{உடனுக்குடன்} (uṭaṉukkuṭaṉ) — then and there; promptly}

\label{lesson:TA-C317-utanukkutan}



\begin{quote}
Before the new one: say the Tamil for again and again, then the Tamil for continuously.
\end{quote}

\subsection*{You'll want to know: \ta{உடனுக்குடன்}}
\textbf{\ta{உடனுக்குடன்}} — \emph{uṭaṉukkuṭaṉ} — \textquotedblleft{}then and there; promptly\textquotedblright{}.

\begin{quote}
\ta{நிருபர்கள்} \ta{செய்திகளை} \ta{உடனுக்குடன்} \ta{அனுப்பினார்கள்}.
\end{quote}

\begin{quote}
\emph{niruparkaḷ ceytikaḷai uṭaṉukkuṭaṉ aṉuppiṉārkaḷ.}
\end{quote}

\begin{quote}
The reporters sent the news then and there.
\end{quote}

\ta{உடன்} doubled with -\ta{க்கு}: as soon as anything happens.

\begin{cousinweb}
\textbf{\ta{உடனுக்குடன்}} (\emph{uṭaṉukkuṭaṉ}) doubles native \ta{உடன்}, \textquotedblleft{}together with, at once\textquotedblright{} (DEDR 945).
\end{cousinweb}

\begin{grammarlens}[title={What you've built}]
One more word for telling what happened.
\end{grammarlens}

\subsection*{Your turn}
\begin{itemize}
\raggedright
  \item \emph{Say it:} \emph{uṭaṉukkuṭaṉ}
  \item \emph{Say it:} \emph{uṭaṉukkuṭaṉ}, once more
  \item \emph{Say it:} \emph{niruparkaḷ ceytikaḷai uṭaṉukkuṭaṉ aṉuppiṉārkaḷ.} — The reporters sent the news then and there.
  \item \emph{Recall:} say the Tamil for again and again, then the Tamil for continuously, then say \emph{uṭaṉukkuṭaṉ} again
\end{itemize}

\begin{tcolorbox}[breakable,colback=teal!4,colframe=teal!35!black,title={Before you move on}]
What does \ta{உடனுக்குடன்} mean? (\textquotedblleft{}Then and there; promptly\textquotedblright{}.) Say it once more.
\end{tcolorbox}
\section[\texorpdfstring{\ta{அடுத்தடுத்து} (aṭuttaṭuttu)}{அடுத்தடுத்து (aṭuttaṭuttu)}]{\ta{அடுத்தடுத்து} (aṭuttaṭuttu) — one after another}

\label{lesson:TA-C317-atuttatuttu}



\begin{quote}
Before the new one: say the Tamil for continuously, then the Tamil for then and there.
\end{quote}

\subsection*{You'll want to know: \ta{அடுத்தடுத்து}}
\textbf{\ta{அடுத்தடுத்து}} — \emph{aṭuttaṭuttu} — \textquotedblleft{}one after another\textquotedblright{}.

\begin{quote}
\ta{நேற்று} \ta{அடுத்தடுத்து} \ta{மூன்று} \ta{விபத்துகள்} \ta{நடந்தன}.
\end{quote}

\begin{quote}
\emph{nēṟṟu aṭuttaṭuttu mūṉṟu vipattukaḷ naṭantaṉa.}
\end{quote}

\begin{quote}
Yesterday three accidents happened one after another.
\end{quote}

\ta{அடுத்து}, next, doubled: each one right after the last.

\begin{cousinweb}
\textbf{\ta{அடுத்தடுத்து}} (\emph{aṭuttaṭuttu}) doubles \ta{அடுத்து}, \textquotedblleft{}next\textquotedblright{}, from native \ta{அடு}, \textquotedblleft{}to be next, near\textquotedblright{} (DEDR 79).
\end{cousinweb}

\begin{grammarlens}[title={What you've built}]
One more word for telling what happened.
\end{grammarlens}

\subsection*{Your turn}
\begin{itemize}
\raggedright
  \item \emph{Say it:} \emph{aṭuttaṭuttu}
  \item \emph{Say it:} \emph{aṭuttaṭuttu}, once more
  \item \emph{Say it:} \emph{nēṟṟu aṭuttaṭuttu mūṉṟu vipattukaḷ naṭantaṉa.} — Yesterday three accidents happened one after another.
  \item \emph{Recall:} say the Tamil for continuously, then the Tamil for then and there, then say \emph{aṭuttaṭuttu} again
\end{itemize}

\begin{tcolorbox}[breakable,colback=teal!4,colframe=teal!35!black,title={Before you move on}]
What does \ta{அடுத்தடுத்து} mean? (\textquotedblleft{}One after another\textquotedblright{}.) Say it once more.
\end{tcolorbox}
\section[(dialogue)]{Looking back: from Pictures and Posters to Again and Again}

\label{lesson:TA-R317-review-310-317}



\begin{quote}
Say the Tamil for a cartoon and a poster.
\end{quote}

\subsection*{Your turn}
Say these aloud:

\begin{itemize}
\raggedright
  \item a cartoon, a poster, a documentary, to portray, the media
  \item a slogan, a hunger strike, to rise up, an uprising, a proclamation
  \item in detail, directly, to describe vividly, to recite, oral tradition
  \item meanwhile, immediately, gradually, at long last, earlier
  \item a moment, to recall, over time, by then, following this
  \item to take place, to arise, an event, at present, once in a while
  \item later on, of that day, previous, after that, at the same time
  \item for the first time, again and again, continuously, then and there, one after another
\end{itemize}

\begin{tcolorbox}[breakable,colback=teal!4,colframe=teal!35!black,title={Before you move on}]
A cartoon? (\textbf{\ta{கேலிச்சித்திரம்}}, \emph{kēliccittiram}.) A poster? (\textbf{\ta{சுவரொட்டி}}, \emph{cuvaroṭṭi}.) A documentary? (\textbf{\ta{ஆவணப்படம்}}, \emph{āvaṇappaṭam}.) To portray? (\textbf{\ta{சித்தரி}}, \emph{cittari}.) The media? (\textbf{\ta{ஊடகம்}}, \emph{ūṭakam}.) A slogan? (\textbf{\ta{முழக்கம்}}, \emph{muḻakkam}.) A hunger strike? (\textbf{\ta{உண்ணாவிரதம்}}, \emph{uṇṇāviratam}.) To rise up? (\textbf{\ta{கிளர்ந்தெழு}}, \emph{kiḷarnteḻu}.) An uprising? (\textbf{\ta{எழுச்சி}}, \emph{eḻucci}.) A proclamation? (\textbf{\ta{பிரகடனம்}}, \emph{pirakaṭaṉam}.) In detail? (\textbf{\ta{விரிவாக}}, \emph{virivāka}.) Directly? (\textbf{\ta{நேரடியாக}}, \emph{nēraṭiyāka}.) To describe vividly? (\textbf{\ta{வர்ணி}}, \emph{varṇi}.) To recite? (\textbf{\ta{ஒப்பி}}, \emph{oppi}.) Oral tradition? (\textbf{\ta{வாய்மொழி}}, \emph{vāymoḻi}.) Meanwhile? (\textbf{\ta{இதற்கிடையில்}}, \emph{itaṟkiṭaiyil}.) Immediately? (\textbf{\ta{உடனடியாக}}, \emph{uṭaṉaṭiyāka}.) Gradually? (\textbf{\ta{படிப்படியாக}}, \emph{paṭippaṭiyāka}.) At long last? (\textbf{\ta{ஒருவழியாக}}, \emph{oruvaḻiyāka}.) Earlier? (\textbf{\ta{முன்னதாக}}, \emph{muṉṉatāka}.) A moment? (\textbf{\ta{தருணம்}}, \emph{taruṇam}.) To recall? (\textbf{\ta{நினைவுகூர்}}, \emph{niṉaivukūr}.) Over time? (\textbf{\ta{நாளடைவில்}}, \emph{nāḷaṭaivil}.) By then? (\textbf{\ta{அதற்குள்}}, \emph{ataṟkuḷ}.) Following this? (\textbf{\ta{இதையடுத்து}}, \emph{itaiyaṭuttu}.) To take place? (\textbf{\ta{நடைபெறு}}, \emph{naṭaipeṟu}.) To arise? (\textbf{\ta{ஏற்படு}}, \emph{ēṟpaṭu}.) An event? (\textbf{\ta{நிகழ்வு}}, \emph{nikaḻvu}.) At present? (\textbf{\ta{தற்போது}}, \emph{taṟpōtu}.) Once in a while? (\textbf{\ta{எப்போதாவது}}, \emph{eppōtāvatu}.) Later on? (\textbf{\ta{பின்னர்}}, \emph{piṉṉar}.) Of that day? (\textbf{\ta{அன்றைய}}, \emph{aṉṟaiya}.) Previous? (\textbf{\ta{முந்தைய}}, \emph{muntaiya}.) After that? (\textbf{\ta{அதன்பின்}}, \emph{ataṉpiṉ}.) At the same time? (\textbf{\ta{அதே} \ta{சமயம்}}, \emph{atē camayam}.) For the first time? (\textbf{\ta{முதன்முறையாக}}, \emph{mutaṉmuṟaiyāka}.) Again and again? (\textbf{\ta{திரும்பத்திரும்ப}}, \emph{tirumpattirumpa}.) Continuously? (\textbf{\ta{தொடர்ச்சியாக}}, \emph{toṭarcciyāka}.) Then and there? (\textbf{\ta{உடனுக்குடன்}}, \emph{uṭaṉukkuṭaṉ}.) One after another? (\textbf{\ta{அடுத்தடுத்து}}, \emph{aṭuttaṭuttu}.)
\end{tcolorbox}
